\section{A sharp matrix classification}\label{sec:classification}
In this section only, let \(A\in\Z^{r\times n}\) be an arbitrary full-row-rank integer matrix, \(1\le r\le n\); it is not initially assumed unimodular. Write \(F^A_b\) for its integer-box fiber enumerator. All boxes have finite integral endpoints and may fix coordinates. We say that \(A\) is \emph{equimodular} if all its nonzero maximal minors have the same absolute value \(D>0\). This is classical terminology; some sources call this class unimodular without requiring \(D=1\)~\cite{BS,CGR}.

The properties below are \emph{universal}: each is quantified over every such fixed box and every triple \(b_-,b_0,b_+\in\Z^r\) with \(b_-+b_+=2b_0\). Let \(\mathsf C_+\) mean coefficientwise nonnegativity of
\[
 (F^A_{b_0})^2-F^A_{b_-}F^A_{b_+}
\]
whenever the three \emph{integer-point sets} are nonempty. Define \(\mathsf W_+\) by requiring the numerical inequality at every positive field, and \(\mathsf N_+\) by requiring it at the field \(\mathbf1\). A merely nonempty real fiber is not this promise.

\begin{theorem}[Positive-fiber classification]\label{thm:classification}
For every full-row-rank integer \(A\), the following are equivalent:
\[
 \mathsf C_+,\qquad\mathsf W_+,\qquad\mathsf N_+,
 \qquad A\text{ is equimodular}.
\]
If \(A\) is not equimodular, a primitive minimal kernel circuit gives a box with at most \(r+1\) nonfixed coordinates and three distinct midpoint right-hand sides with counts \((1,1,2)\). Their coefficient defect is one negative monomial, so the same triple violates the inequality at \emph{every} positive field.
\end{theorem}

The equivalence of the three inequalities is at the level of these universal properties. A numerical inequality for one arbitrary triple is not asserted to imply coefficientwise nonnegativity. The sufficiency below is a classical normalization followed by the external inequality~\eqref{eq:input}. The possible additional observation is the necessity with an explicit \emph{all-positive} count witness, rather than a central lattice hole alone.

\begin{lemma}[Classical circuit criterion]\label{lem:circuits}
A full-row-rank integer matrix is equimodular if and only if every primitive minimal kernel circuit has all its nonzero coefficients in \(\{1,-1\}\).
\end{lemma}
\begin{proof}
The statement is recorded explicitly in~\cite[Definition 1.1(2)--(3)]{BS}; we recall the elementary argument to fix conventions. A circuit is a nonzero integer kernel vector of minimal support, with relatively prime nonzero coefficients. Its support \(S\) has size at most \(r+1\), and the restricted real kernel is one-dimensional.

If all nonzero maximal minors have absolute value \(D\), extend \(S\setminus\{j\}\) to a column basis. The resulting fundamental circuit has support \(S\); Cramer's rule expresses its nonzero coefficient ratios as ratios of nonzero maximal minors. Their absolute values are one, and primitive scaling gives coefficients \(\pm1\). A zero column has a singleton circuit with primitive coefficient one.

Conversely, connect two column bases \(I,J\) by exchanges. For \(j\in J\setminus I\), its fundamental circuit in \(I\cup\{j\}\) contains an \(i\in I\setminus J\), since otherwise \(J\) would contain a dependence. Exchanging \(i\) for \(j\) increases the intersection with \(J\). The determinant ratio of the two bases is the corresponding circuit-coefficient ratio, hence has absolute value one. Iteration proves equimodularity. This is a classical characterization, not a new recognition theorem.
\end{proof}

\begin{lemma}[A positive circuit certificate]\label{lem:positive-circuit}
Let \(c\) be a primitive minimal integer kernel circuit of \(A\), oriented so that \(c_i=a\ge2\). Put \(p=c^+\), \(q=c^-\), and use the box \(0\le x\le |c|\), fixing all coordinates outside its support to zero. For the column \(A_i\), set
\[
 R_k=Ap-kA_i\quad(k=0,1,2).
\]
Then
\[
 F^A_{R_0}=z^p+z^q,\quad
 F^A_{R_1}=z^{p-e_i},\quad
 F^A_{R_2}=z^{p-2e_i},
\]
and consequently
\begin{equation}\label{eq:negative-monomial}
 (F^A_{R_1})^2-F^A_{R_2}F^A_{R_0}
   =-z^{|c|-2e_i}.
\end{equation}
\end{lemma}
\begin{proof}
The point \(x_k=p-ke_i\) lies in the box for each \(k\). Every other integer solution on the circuit support is \(x_k+tc\), with \(t\in\Z\): the real kernel is \(\R c\), and a Bezout combination of the coefficients of the primitive \(c\) shows that \(tc\in\Z^n\) forces \(t\in\Z\).

The \(i\)-th coordinate imposes \(-1+k/a\le t\le k/a\). There is another support coordinate \(j\ne i\), since a primitive singleton circuit cannot have coefficient of magnitude at least two. Its bound imposes \(-1\le t\le0\), regardless of the sign of \(c_j\). Together the bounds give
\[
 -1+k/a\le t\le0.
\]
For \(k=0\) the admissible integers are \(-1,0\), yielding \(q,p\); for \(k=1,2\), the only one is zero. This includes \(a=2\), uniform-sign circuits and circuits with no coefficient of magnitude one. The displayed enumerators and defect follow. The exponent \(|c|-2e_i\) is nonnegative. Finally \(A_i\ne0\), since a zero column would be a smaller singleton dependence, so the three right-hand sides are distinct.
\end{proof}

\begin{proof}[Proof of Theorem~\ref{thm:classification}]
The implications \(\mathsf C_+\Rightarrow\mathsf W_+\Rightarrow\mathsf N_+\) are immediate. If \(A\) is not equimodular, Lemma~\ref{lem:circuits} supplies the circuit in Lemma~\ref{lem:positive-circuit}. The midpoint triple \((R_2,R_1,R_0)\) has positive counts \((1,1,2)\), contradicting \(\mathsf N_+\). Its support has size at most \(r+1\), and~\eqref{eq:negative-monomial} gives the stronger field-independent failure.

For sufficiency, choose a nonsingular column basis \(B\) of an equimodular \(A\), with \(|\det B|=D\). Cramer's rule gives \(C=B^{-1}A\in\Z^{r\times n}\), with entries \(0,\pm1\) and maximal minors \(0,\pm1\). Since \(C\) contains an identity column basis,
\begin{equation}\label{eq:column-lattice}
 A\Z^n=B\Z^r.
\end{equation}
Each nonempty integer fiber has \(b_s=Ax_s\) for an integer \(x_s\), and therefore \(B^{-1}b_s=Cx_s\in\Z^r\). The transformed integral right-hand sides retain their midpoint relation. The transformation changes only row equations, preserving the exact real and integer fibers, box and original coordinate fields. Apply Leake--Mohammadi Yekta's coefficientwise Theorem 3.1 to \(C\). This proves \(\mathsf C_+\). The normalization and lattice identity are classical~\cite[Theorem 3]{CGR}; \(B^{-1}\) itself need not be integral.
\end{proof}

\begin{theorem}[Ambient lattice and parity]\label{thm:odd}
Remove the promise that all three integer fibers are nonempty. Then each universal inequality holds if and only if \(A\) is equimodular and its common nonzero maximal-minor magnitude \(D\) is odd. For equimodular even \(D\), a binary box with at most \(r\) nonfixed coordinates gives midpoint counts \((1,0,1)\).
\end{theorem}
\begin{proof}
Necessity of equimodularity follows from Theorem~\ref{thm:classification}. If \(D\) is even, a column basis \(B\) is singular modulo two. Choose nonzero \(\alpha\in\{0,1\}^r\) with \(B\alpha\) even. On those basis coordinates take the box \([0,\alpha]\), fixing the others to zero. Invertibility of \(B\) makes the extreme fibers at \(0,B\alpha\) the singletons \(0,\alpha\). The integral midpoint right-hand side \(B\alpha/2\) has the unique real solution \(\alpha/2\), which is not integral. This gives the claimed failure.

Now suppose \(D\) is odd. If an endpoint enumerator is zero, the coefficientwise inequality is immediate. Otherwise \(b_-,b_+\in B\Z^r\) by~\eqref{eq:column-lattice}. For \(v=B^{-1}b_0\), both \(2v\) and \(Dv\) are integral, the latter by the adjugate formula. Bezout's identity for \(2,D\) makes \(v\) integral. All transformed right-hand sides are integral, so the same external theorem for \(C\) proves coefficientwise nonnegativity, including the initially possibly empty central fiber. Numerical versions follow as before.
\end{proof}

\begin{corollary}[Parity-halving formulation]\label{cor:halving}
The unconditional universal inequalities in Theorem~\ref{thm:odd} are equivalent to: for every \(\alpha\in\{0,1\}^n\) with \(A\alpha\in2\Z^r\), there is an integer \(\beta\) with \(0\le\beta\le\alpha\) and \(2A\beta=A\alpha\).
\end{corollary}
\begin{proof}
The unconditional inequality applied to \([0,\alpha]\) forces a nonempty integer central fiber because both endpoints are nonempty. Conversely, for a bad primitive circuit \(c\), let \(\alpha\) be its coordinatewise residue modulo two. It is nonzero and \(A\alpha\) is even. A hypothetical \(\beta\) gives the nonzero kernel vector \(d=2\beta-\alpha\) with support contained in that of \(c\) and nonzero entries \(\pm1\). Minimal support and the one-dimensional kernel force proportionality; primitivity would then make every nonzero coefficient of \(c\) have magnitude one, a contradiction. Thus \(A\) is equimodular. The even-basis construction above excludes even \(D\), and Theorem~\ref{thm:odd} finishes the proof.
\end{proof}
This feasibility formulation is close to the existing proof, not a separate new geometric class. For example write \(\ell+u=2h+\alpha\), with \(h\in\Z^n\) and binary \(\alpha\). If \(A(\ell+u)\) is even, the criterion supplies \(h+\beta\) in the box with image \(A(\ell+u)/2\); this is the central existence step used in~\cite[Lemma 3.3]{LMY}.

\paragraph{Equality and algorithm transfer.}
For equimodular \(A\) and three nonempty integer fibers, the same normalization transfers Theorem~\ref{thm:main} and Theorem~\ref{thm:alg} below without changing polytopes or fields. Their equality characterization and \(6n\) linear-optimization decision therefore apply on the column lattice. This is an application of the classical row normalization; it does not provide a new equimodularity recognition algorithm or remove regular kernel geometry.

\paragraph{A limit, not a Jensen solution.}
The odd-index extension is specific to midpoints. For \(A=[3]\), the fibers at right-hand sides \(0,3\) in \([0,1]\) have count one, while their mean with weights \(2/3,1/3\) is the integral right-hand side \(1\), whose count is zero. The same counts hold for unbounded nonnegative fibers of this matrix. This matrix is not unimodular in the \(D=1\) convention of the open question; it does not answer~\cite[Question 1.7]{LMY}.
